\section{研究背景}\label{背景}
% !!自己開示と，開示に対する受容が大事 → しかし人間関係の希薄化が課題 というのにきれいにつなげたい．

現代社会において，人々の幸福な生活の実現には親密な人間関係の構築が重要な要素の1つである．社会心理学の分野において，親密な関係を築く上で自己開示（self-disclosure）が重要なプロセスであることが明らかにされている\cite{collins_1994}．この自己開示の研究は，Jourardをはじめとして1958年から心理学の分野で行われており，その知見は現代においても重要性が増している．

しかし，和田\cite{和田実_1990_青年の対人関係の変容}は，日本の若者が自分の趣味などについて表面的には自己開示を行っているように見えるものの，自分の内面性に触れるような深いレベルの自己開示は行えていないと指摘している．土井\cite{土井隆義_2004_個性}も同様に，ある程度までは他者と仲良くなることができても，さらに一歩進んだ関係にまでは誰とも親密な関係を発展させることができない若者が増えていると述べている．

さらに近年，SNSの普及による表面的なコミュニケーションの増加や，新型コロナウイルス感染症（COVID-19）の影響による対面での交流機会の減少により，若者の人間関係の希薄化が社会的な課題となっている．特に，大学生において，内面的な友人関係を避ける傾向\cite{岡田努_2007}や，他者からの評価を気にして友人関係の深まりを避ける関係回避特徴が指摘されている\cite{岡田努_2010_青年期の友人関係と自己}．このような人間関係の希薄化は，現代社会における重要な課題である．

